%=============================================================================!
% 18.4  KERNELS AND GAUSSIAN PROCESSES                                        !
%=============================================================================!
\section{Kernels and Gaussian processes}
\label{sec:lo-gp}

A forecaster who knows the temperature at a few weather stations
estimates it in a village between them by leaning on the nearest
stations most, since temperatures a few kilometres apart are alike, and
says how sure they are: close to a station they are nearly certain, far
from every station they can say little more than what is usual for the
season. A Gaussian process makes this reasoning exact. It needs no basis
functions chosen in advance, only a statement of how alike the values at
two points are expected to be, and it returns, at every point, an
estimate and an error bar.

\subsection*{A random function}

Before any reading, the unknown curve $f(x)$ is treated as random: its
values at any set of points are jointly Gaussian with mean zero, and the
values at two points $x$ and $x'$ have the covariance
\begin{equation}
   k(x, x') = h^2\exp\Bigl[-\frac{(x - x')^2}{2\ell^2}\Bigr] ,
   \label{eq:lo-kernel}
\end{equation}
the \emph{kernel}. Each value has the variance $k(x, x) = h^2$, a spread
$h$ about zero; values at points much closer than the length scale $\ell$
are nearly equal, their covariance almost $h^2$, and values much further
apart are nearly independent. The kernel is the forecaster's belief that
nearby temperatures are alike, written as a number. Readings $y_i =
f(x_i) + \epsilon_i$ with independent noise $\epsilon_i$ of spread
$\sigma$ then have the covariances $K_{ij} = k(x_i, x_j) +
\sigma^2\delta_{ij}$, the noise adding to each reading's own variance
only. The value $f(x_\ast)$ at a new point and the readings are jointly
Gaussian, with the covariances $k(x_\ast, x_i)$, collected in the vector
$\vb{k}_\ast$. The question is what the readings say about $f(x_\ast)$:
the Gaussian conditioned on them.

\begin{toolbox}[A Gaussian conditioned on another]
Let $a$ and $b$ be jointly Gaussian with mean zero, variances
$\sigma_a^2$ and $\sigma_b^2$ and covariance $c = \ens{ab}$. The number
$e = a - (c/\sigma_b^2)\,b$ has the covariance with $b$
\begin{equation*}
   \ens{eb} = \ens{ab} - \frac{c}{\sigma_b^2}\ens{b^2}
   = c - \frac{c}{\sigma_b^2}\sigma_b^2 = 0 .
\end{equation*}
A linear combination of jointly Gaussian numbers is Gaussian, and two
jointly Gaussian numbers with zero covariance are independent, since
their joint density, $\exp(-e^2/2\sigma_e^2 - b^2/2\sigma_b^2)$ up to a
constant, is the product of the two densities (\cref{sec:sm-probability}).
So knowing $b$ says nothing about $e$, and $a = (c/\sigma_b^2)b + e$ has,
given $b$, the mean $(c/\sigma_b^2)\,b$ and the variance of $e$,
\begin{align*}
   \sigma_e^2 &= \ens{a^2} - 2\frac{c}{\sigma_b^2}\ens{ab}
      + \frac{c^2}{\sigma_b^4}\ens{b^2}
      && \text{expanding the square of $e$}\\
   &= \sigma_a^2 - \frac{2c^2}{\sigma_b^2} + \frac{c^2}{\sigma_b^2}
   = \sigma_a^2 - \frac{c^2}{\sigma_b^2}
      && \text{collecting the terms in $c^2$.}
\end{align*}
The same steps with $b$ a vector of readings, of covariance matrix
$K$, and $\vb{c} = \ens{a\vb{b}}$ a vector, give the mean $\vb{c}^{\mathsf
T}K^{-1}\vb{b}$ and the variance $\sigma_a^2 - \vb{c}^{\mathsf
T}K^{-1}\vb{c}$, the division by $\sigma_b^2$ becoming the inverse matrix
(\cref{sec:ro-remove}). Knowing $b$ moves the mean and always lowers
the variance.
\end{toolbox}

\noindent With $a = f(x_\ast)$, $\vb{b} = \vb{y}$ and $\vb{c} =
\vb{k}_\ast$, the readings give $f(x_\ast)$ the mean and variance
\begin{equation}
   \bar f(x_\ast) = \vb{k}_\ast^{\mathsf T}K^{-1}\vb{y} , \qquad
   \operatorname{var} f(x_\ast) = h^2 - \vb{k}_\ast^{\mathsf T}K^{-1}
   \vb{k}_\ast .
   \label{eq:lo-gp}
\end{equation}
The mean is a weighted sum of the readings, as the forecaster's estimate
is, the weights $K^{-1}\vb{k}_\ast$ larger for readings near $x_\ast$.
Close to a reading the variance falls to about the noise's $\sigma^2$;
far from all of them $\vb{k}_\ast$ vanishes and the variance returns to
the prior's $h^2$. This is the posterior of the weights of
\cref{sec:lo-ridge} in another form: a Gaussian process is ridge
regression on as many basis functions as the kernel implies, possibly
infinitely many, solved through the $n\times n$ matrix $K$ of the
readings rather than the $p\times p$ matrix $X^{\mathsf T}X$ of the
weights\cite{Rasmussen2006}. \texttt{mdlab.learn.GaussianProcess} solves
it with the factor $K = LL^{\mathsf T}$, $L$ lower triangular, by which a
system with $K$ becomes two systems with triangles, each solved one
unknown at a time.

\subsection*{The pair energy from eight readings}

Eight readings of argon's pair energy, drawn as in \cref{sec:lo-basis},
are fitted with $h = \SI{10}{\milli\electronvolt}$, about the depth of the
well, and the known noise of \SI{0.2}{\milli\electronvolt}
(\cref{fig:lo-gp}). The length scale still has to be chosen, and the
readings themselves choose it through their \emph{marginal likelihood},
the density of the readings under the prior. Since $\vb{y}$ is Gaussian
with covariance $K$, its logarithm is
\begin{equation*}
   \ln p(\vb{y}) = -\tfrac12\vb{y}^{\mathsf T}K^{-1}\vb{y}
   - \tfrac12\ln\det K - \tfrac n2\ln2\pi ,
\end{equation*}
the multivariable form of the Gaussian of \cref{sec:sm-maxwell}. The
first term rewards a prior under which the readings are typical; the
second penalises a prior so loose, with $\ell$ so short, that it would
have found any readings typical. Their balance peaks at $\ell =
\SI{0.70}{\angstrom}$ (\cref{fig:lo-gp}c). There the mean follows the
readings, and its standard deviation is at most
\SI{0.2}{\milli\electronvolt} at the readings and
\SI{1.1}{\milli\electronvolt} midway across the widest gap, from 3.36 to
\SI{4.13}{\angstrom}. A length a quarter as long makes the curve sag
towards zero between readings, with an honest but useless band; one four
times as long makes it too stiff, with a band too narrow
(\cref{fig:lo-gp}b).

The band must also be right. Over the test, \num{0.75} of the noise-free
values lie within twice the standard deviation at $\ell =
\SI{0.70}{\angstrom}$, against the 0.95 that a correct Gaussian band
holds, and the misses gather at the steep wall below \SI{3.8}{\angstrom}
and at the bottom of the well, where no reading falls. A kernel with one
length scale cannot be both smooth enough for the slow tail and supple
enough for the wall, and the error bar is only as good as the kernel's
belief. At a quarter of the length the band holds 0.99 of the test, and
at four times only 0.15. Gaussian processes on atoms, with kernels built
on descriptions of each atom's surroundings, are the Gaussian
Approximation Potentials\cite{Bartok2010}, among the first learned
potentials; their error bar is what lets a learned potential say where
it needs another first-principles calculation, the subject of
planned Chapter~23.

\begin{figure}[htb]
\centering
\includegraphics{ch18_learning/gp}
\caption{A Gaussian process through eight readings of argon's pair
energy (dots, noise \SI{0.2}{\milli\electronvolt}), with $h =
\SI{10}{\milli\electronvolt}$, against the pair energy (grey dashed). (a)
The posterior mean and twice its standard deviation (band) for the length
scale of greatest marginal likelihood. (b) The same for a length a
quarter of that and four times it. (c) The log marginal likelihood of the
readings against the length scale, with its greatest value (dashed).}
\label{fig:lo-gp}
\end{figure}

\begin{keyidea}
A Gaussian process treats the unknown curve as a random function whose
values covary by a kernel; conditioning on the readings gives the mean
$\vb{k}_\ast^{\mathsf T}K^{-1}\vb{y}$ and the variance $h^2 -
\vb{k}_\ast^{\mathsf T}K^{-1}\vb{k}_\ast$ at every point, an estimate with
an error bar that grows away from the readings. The marginal likelihood
chooses the kernel's length. The error bar is only as reliable as the
kernel, and checking it against held-out values is part of the fit.
\end{keyidea}

\notebook{18}{move the readings and the length scale and watch the band}

\begin{selfcheck}
A Gaussian process is conditioned on one reading $y_1$ at $x_1$, with no
noise. What are its mean and variance at $x_1$, and at a point very far
from $x_1$?
\end{selfcheck}
